\documentclass[12pt]{amsart}
\pagestyle{empty}
\usepackage[hmargin=1in,vmargin=1in]{geometry}
\usepackage{url,  mathrsfs}
\usepackage{hyperref}
\usepackage{fancyhdr} 
\usepackage[T1]{fontenc}


\usepackage{url,  mathrsfs}
\usepackage{hyperref}
\usepackage{amsmath}
\usepackage{graphicx,  epstopdf}
\usepackage{subfig}
\usepackage{color}
\usepackage{bbm}
\usepackage{subfloat}
\usepackage[draft]{fixme}
\usepackage{bm} 
\usepackage{bbm} 
\usepackage{epigraph}
\usepackage{endnotes}
\usepackage{ifpdf}
\usepackage{array}

\usepackage{multicol}
\usepackage{multirow}
\usepackage{graphicx}
\usepackage{tikz}

\newcommand{\A}{\mathbb{A}}
\newcommand{\Q}{\mathbb{Q}}
\newcommand{\N}{\mathbb{N}}
\newcommand{\Z}{\mathbb{Z}}
\newcommand{\R}{\mathbb{R}}
\newcommand{\F}{\mathbb{F}}
\newcommand{\C}{\mathbb{C}}
\renewcommand{\P}{\mathbb{P}}
\newcommand{\I}{\mathcal{I}}
\newcommand{\V}{\mathcal{V}}
\newcommand{\cH}{\mathcal{H}}
\newcommand{\ol}{\overline}
\newcommand{\J}{\mathcal{J}}
\newcommand{\X}{\mathcal{X}}
\theoremstyle{definition}
\newtheorem{prob}{Problem}
\newtheorem{extraProb}{Extra Problem}

\DeclareMathOperator{\conv}{conv}
\DeclareMathOperator{\cone}{cone}
\DeclareMathOperator{\ext}{ext}


\usepackage{listings,xcolor}
\lstset{language=Mathematica}
\lstset{basicstyle={\sffamily\footnotesize},
  numbers=left,
  numberstyle=\tiny\color{gray},
  numbersep=5pt,
  breaklines=true,
  captionpos={t},
  frame={lines},
  rulecolor=\color{black},
  framerule=0.5pt,
  columns=flexible,
  tabsize=2
}

\begin{document}
 \thispagestyle{empty}
\begin{center} {\large \textbf{Math 380 -- Homework  2}} \smallskip \\ 
Due on Thursday, April 16\ \\ \ \\
 \end{center}

You are welcome to talk with other students in the class about problems but should write up solutions on your own. 
Solutions can be handwritten or typed but need to be legible and submitted via Gradescope by the end of the day 
on Thursday. You should justify all your answers in order to receive full credit.  
\bigskip 

\begin{flushleft}

{\bf Good practice problems} (not to be turned in): \\

CLO Section 1.4, Exercises 2, 3, 6, 8, 16\\

CLO Section 1.5, Exercises 1, 11, 13, 14, 16\\
CLO Section 2.1, Exercises 1, 3, 4\\
\bigskip 




\begin{prob} CLO \S1.4, Exercise 15 \smallskip \\
\begin{enumerate}
    \item[a.] For any subset $S \subseteq k^n$, the set $\mathbf I(S) = \{f \in k[x_1,\dots,x_n] \mid \forall a \in S, f(a) = 0 \}$ is an ideal. First, $0 \in \mathbf I(S)$ as the zero polynomial vanishes on all of $k^n$ and thus also vanishes on $S$. Let $f,g \in \mathbf I(S)$ be any polynomials, and consider $f+g$ evaluated at any point $s \in S$. Because $f$ and $g$ vanish on $S$, $(f+g)(s)=f(s)+g(s)=0+0=0$. Thus, $f+g \in \mathbf I(S)$. Now, let $f \in \mathbf I(S)$ and let $h\in k[x_1,\dots,x_n]$, then consider $h \cdot f$ evaluated at $s \in S$. Since $f$ vanishes on $S$, $(h \cdot f)(s) = h(s)\cdot f(s)=h(s)\cdot 0 = 0$, so $h\cdot f \in I$. Thus, the set $\mathbf I(S)$ is an ideal.
    \item[b.] Let $X = \{(a,a) \in \mathbb R^2 \mid a \ne 1 \}$. We would like to determine the ideal $\mathbf I(X)$. First, note that $\mathbf I(X) = \mathbf I(V(y-x))$: we have $\mathbf I(X) \subseteq \mathbf I(V(y-x))$ since by Exercise 8 of \S 1.2, any polynomial $f$ that vanishes on $X$ also vanishes on $(1,1)$, so $f$ vanishes on the entire $y=x$ line. The other direction is true since $X \subseteq V(y-x)$, and so $\mathbf I(X) \supseteq \mathbf I(V(y-x))$. Any polynomial that vanishes on all of $V(y-x)$ must also vanish on a subset $X$. By Exercise 10 of \S 1.4, $\mathbf I(V(y-x)) = \langle y-x\rangle$, and so we have found $\mathbf I(X) = \langle y-x\rangle$. 
    
    Further proof of Exercise 10 for fun: the containment that $\langle y-x\rangle \subseteq \mathbf I(V(y-x))$ is trivial, since $y-x$ vanishes on the the line $y=x$. For the other containment, let $f \in \mathbf I(\mathbf V(y-x))$ be any polynomial that vanishes on $\mathbf V(y-x)$. Using the binomial theorem, we can write any monomial in $k[x,y]$ as \begin{align*}x^ny^m &= x^n(x+(y-x))^m
    \\ &= x^n\left(x^m +(y-x)\sum_{i=1}^m \binom{m}{i}(y-x)^{i-1}x^{m-i}\right)
    \\ &= x^{n+m} + (y-x)\cdot h(x).
    \end{align*}
    Because $f$ is a $\mathbb R$-linear combination of these monomials, we can write $f(x,y) = p(x)+(y-x)\cdot h(x)$ where $p(x), h(x) \in \mathbb R[x]$. Because $f$ vanishes everywhere on the $y=x$ line, we have that $\forall x \in \mathbb R, f(x,x)=p(x)+0=0$. Thus, $p(x)$ is a univariate polynomial with infinite roots, so it is the zero polynomial. Thus, $(y-x)$ divides $f$, and $f \in \langle y-x\rangle$.
    %and choose any $a \in k$. Then, consider the univariate polynomial $p(y)=f(a,y)$. \textbf{TODO: Take care for case when $\deg p < 1$} We can divide $p(y)=q(y)\cdot (y-a)+r(y)$, and by the degree bound we have that $r(y) \in k$ is a constant polynomial. Thus, evaluating $p(a)=0=r(y)$, we see that $y-a$ divides $p(y)$.  Surely this implies that $y-x$ divides $f$, but I dont know how to do it.
    \item[c.] Let $\mathbb Z^n$ be the set of all points in $\mathbb C^n$ with integer coordinates. Then $\mathbf I(\mathbb Z^n) = \langle 0 \rangle$. Consider any polynomial $f \in \mathbb C[x]$ that vanishes on all of $\mathbb Z^n$. By the result of Exercise 6 of \S 1.1, $f = 0$. The proof is the same as Proposition 5, inducting on $n$: for the $n=1$ case, $f$ is a univariate polynomial with infinite roots and so it is the zero polynomial. For $n > 1$, write \[f = \sum_{i=0}^{\deg f} g_i(x_1, \dots, x_{n-1}) x_n^i\] for $g_i \in k[x_1,\dots,x_{n-1}]$.
    Let $a \in \mathbb Z^{n-1}$ be any lattice point, and consider the univariate polynomial $p(x_n)= f(a,x_n)= \sum_i g_i(a) x_n^i$. Then for all $x_n \in \mathbb Z$, $p(x_n) = 0$ so by the $n=1$ case we have $p = 0$. Thus, for all $0 \le i \le \deg f$ we have $g_i(a) = 0$. We have shown that all $g_i$ vanish at all points $a \in \mathbb Z^{n-1}$, so by the inductive hypothesis, $g_i = 0$. Thus, $f = 0$. 
\end{enumerate}
\end{prob}
\bigskip 


\begin{prob} CLO \S1.5, Exercise 12 and 15 \smallskip \\
You may assume the results of Exercise 1, 13, and 14 of this section without proof. 

\paragraph{Problem 12} Let \[f = c \prod_{i=1}^l(x-a_i)^{r_i} \in \mathbb C[x]\] be a nonconstant polynomial with $c \ne 0$ and $l \ge 1$. Define the reduced or square-free part of $f$ to be \[f_\text{red} = c \prod_{i=1}^l(x-a_i).\] \begin{enumerate}
    \item[a.] $\mathbf V(f) = \{a_1,\dots,a_l\}$. The containment that $\{a_1, \dots, a_l\}\subseteq \mathbf V(f)$ follows from the definitions, as plugging in any $a_i$ will result in $0$. For the reverse containment, choose any $b \in k$ such that $b \not\in \{a_1,\dots,a_l\}$. Then $f(b) = c\prod (b-a_i)^{r_i}$, and each $b-a_i \ne 0$. However, then $f(b)$ is a product of only nonzero terms, so since $\mathbb C$ is an integral domain, $f(b)\ne 0$. Thus, $b \not\in \{a_1,\dots, a_l\} \implies f(b) \ne 0$, and the contrapositive is our desired containment.
    \item[b.] $\mathbf I(\mathbf V(f)) = \langle f_\text{red}\rangle$. The containment that $\langle f_\text{red}\rangle \subseteq \mathbf I(\mathbf V(f))$ follows from the definitions since $f_\text{red}$ vanishes at every $a_i$ and thus vanishes at every point that $f$ vanishes on. To show the nontrivial containment that $\textbf I(\textbf V(f)) \subseteq \langle f_\text{red} \rangle$, let $g \in \mathbf I(\mathbf V(f))$ be any polynomial that vanishes on $\{a_1,\dots,a_l\}$. We can divide $g = q \cdot f_\text{red} + r$, and because $g$ and $f_\text{red}$ both vanish on all $a_i$, we have $r(a_i) = 0$ for all $1 \le i \le l$. By construction, $\deg r < \deg f_\text{red} = l$, but we have just shown that $r$ is a polynomial with at least $l$ roots! Thus, $r$ is the constant polynomial $0$, and $f_\text{red}$ divides $g$.

    Is there a slicker way to do this using the fact that $\mathbb C[x]$ is a principal ideal domain?
\end{enumerate}

\paragraph{Problem 15} \begin{enumerate}
    \item[a.] By part (c) of \S1.5 Problem 14, given $f = c \prod_{i=1}^l (x-a_i)^{r_i} \in \mathbb C[x]$, we have $\gcd(f,f')=\prod_{i=1}^l (x-a_i)^{r_i-1}$ where $f'$ is the formal derivative of $f$. 
    Then, since we have that each $r_i \ge 1$,
    \begin{align*}
        \frac{f}{\gcd(f,f')} &= \frac{c\prod_{i=1}^l  (x-a_i)^{r_i}}{\prod_{i=1}^l (x-a_i)^{r_i-1}} \\
        &= c\prod_{i=1}^l (x-a_i) = f_\text{red}.
    \end{align*}
    \item[b.] 
    We can calculate the square-free form of $x^{11} - x^{10} + 2 x^8 - 4 x^7 + 3x^5 - 3x^4 + x^3 + 3x^2 -x - 1$ using Mathematica:
    \begin{lstlisting}[language=Mathematica]
        f[x_] := x^11 - x^10 + 2 x^8 - 4 x^7 + 3 x^5 - 3 x^4 + x^3 + 3 x^2 - x - 1
        f[x]/PolynomialGCD[f[x], f'[x]] // Simplify
    \end{lstlisting}
    The answer being $x^5+x^2-x-1$.
 \end{enumerate}
\end{prob}
\bigskip 

\begin{prob} CLO \S2.1, Exercise 5 \smallskip \\

\begin{enumerate}
    \item[a.]  The total number of monomials $x^ay^b \in k[x,y]$ of degree less than or equal to $m$ is given by the triangle numbers, by Mathematica:
    \[\sum_{i=0}^m (i+1) = \frac{(m+2)(m+1)}{2}.\]
    Alternatively, since we need to count the number of pairs $a+b$ where $a+b\le m$, we can rephrase the problem as counting the ways to put $m$ indistinguishable balls into $3$ distinguishable bins: $a$, $b$, or not included. Using stars and bars, have $m$ elements plus $2$ dividers to be chosen, and we find again find \[\binom{m+2}{2} = \frac{(m+2)(m+1)}{2}.\]
    \item[b.] Let $f,g\in k[t]$ be polynomials with degree at most $n$. Fix some $m \in \mathbb N$, and consider the set $S$ of formal ``monomials'' in $f(t)$ and $g(t)$: \[S=\{f(t)^ag(t)^b \mid a,b\in \mathbb Z_{\ge 0}, a+b\le m \}.\] 
    Note that we consider $f^ag^b$ and $f^cg^d$ as distinct elements of $S$ even if $f^ag^b = f^cg^d$ as polynomials in $k[t]$. Denote $P(s)$ as the polynomial in $k[t]$ associated with $s$. For any element $s \in S$, we can see that \[\deg P(s) = \deg (f(t)^a g(t)^b) = a\deg(f) + b\deg(g)\le a n+b n \le nm.\] Thus, $P(S)$ is a subset of polynomials in $t$ with degree at most $nm$, which is a subspace $k[t]_{\deg \le nm}$ of $k[t]$ with $\dim(k[t]_{\deg \le nm}) =nm+1$. Since this is a subspace, we also have $\text{span}\,P(S) \subseteq k[t]_{\deg\le nm}$.

    Note that the definition of $S$ puts it in bijection with pairs $(a,b)$ where $a+b\le m$, so by the counting of part (a) we have that $|S| = \frac{1}{2}(m+2)(m+1)$. Choose something like $m =2n+1$, and doing algebra one can see that \[|S| = (2n+3)(n+1)> 2n^2 + n +1=\dim(k[t]_{\deg\le nm}).\] 
    Since $|S|$ exceeds the dimension of the space that $S$ inhabits, $S$ is linearly dependent (conflating polynomials with the formal monomials). In  fact, any sufficiently large $m$ works since $|S|$ is on the order of $m^2$ while dimension is on the order of $m$.
    \item[c.] Let $C$ be a parametric curve $x=f(t)$, $y=g(t)$ defined by polynomials $f,g\in k[t]$. We can find a nonzero polynomial $F \in k[x,y]$ such that $C \subseteq \mathbf V(F)$. Let $F$ be given by 
    $F = \sum_{i,j} c_{ij} x^i y^j$, and note that we wish for $F$ to have the property that $\forall t \in k, F(f(t),g(t))=0$. Expanding this out, we wish for
    \[F(f(t),g(t)) = \sum_{i,j} c_{ij} f(t)^ig(t)^j=0.\]
    However, by the result of part (b), this list of ``monomials'' $f(t)^i g(t)^j$ must be linearly dependent when we consider large enough total degrees of $F$, so we can find a nontrivial linear combination $c_{ij}$ that sums to zero, proving that our desired nonzero $F$ exists.
    \item[d.] Fix $m \in \mathbb N$ and define a set of formal ``monomials'' for $f,g,h\in k[t,u]$
    \[
    S = \{f(t,u)^ag(t,u)^bh(t,u)^c \mid a+b+ c \le m\}.
    \]
    Let $n$ be the maximum total degree among $f,g,h$, and let $P(s)$ denote the polynomial in $k[t,u]$ associated with any $s \in S$. Then for any $s$,
    \[
    \deg P(s) = a\deg f+b\deg g + c\deg h \le (a+b+c)n\le mn.
    \]
    Thus, $S\subseteq k[t,u]_{\deg \le nm}$ and this is a subspace of $k[t,u]$. However, since this vector space consists of any monomial $t^au^b$ where $a+b\le nm$, by the counting of part (a) we have that \[\dim(k[t,u]_{\deg \le nm})= \frac{1}{2}(nm+1)(nm+2).\] 
    Then, we can use stars and bars to count $|S|$ as the number of ways to put $m$ indistinguishable balls into $4$ distinguishable bins: $a$, $b$, $c$, or not included. Thus, \[|S| = \binom{m+3}{3} = \frac{1}{6}(m+3)(m+2)(m+1).\]
    Since $|S|$ grows on the order of $m^3$ while $\dim(k[t,u]_{\deg\le nm})$ grows on the order of $m^2$, for large enough $m$ we will have that $|S|$ exceeds the dimension of the subspace, so $S$ must be linearly dependent. 

    Consider any parametric surface $C$ given by $x = f(t,u)$, $y = g(t,u)$, and $z = h(t,u)$ for $f,g,h \in k[t,u]$. We wish to show there is a nonzero $F \in k[x,y,z]$ such that $C\subseteq \mathbf V(F)$. To do so, we can write $F(x,y,z) = \sum c_{ijk} x^iy^jz^k$ and show that for all $t,u\in k$, \[
    F(f(t,u),g(t,u),h(t,u)) = \sum c_{ijk} f(t,u)^i g(t,u)^j h(t,u)^k = 0.
    \]
    However, by the result above, the formal monomials $f^i g^j h^k$ must be linearly dependent once we consider large enough degrees of $F$, so we can find a set of nontrivial coefficients $c_{ijk}$ to satisfy the above condition. Thus, $C \subseteq \mathbf V(F)$ for some $F \in k[x,y,z]$.
\end{enumerate}

\emph{Bonus:} Show (c) for rational functions. That is, that any curve in $k^2$ given parametrically by $x=p_1(t)/q_1(t)$ and $y = p_2(t)/q_2(t)$ where $p_1(t), p_2(t), q_1(t), q_2(t)\in k[t]$
is contained in $V(F)$ for some nonzero $F\in k[x,y]$. 
\begin{proof}
    Fix $m \in \mathbb N$ and for $p_1,p_2,q_1,q_2 \in k[t]$ define a set of formal ``monomials''
    \[S = \{ p_1(t)^a p_2(t)^bq_1(t)^{m-a}q_2(t)^{m-b}\mid a+b \le m\}.\]
    This set can be counted in the same way as we did before, so $|S| = \frac{1}{2}(m+2)(m+1)$. Again denote $P(s)$ for the polynomial associated with $s$, and let $n$ be the maximal degree of $p_1,p_2,q_1,q_2$. For any $s \in S$, we see
    \[
    \deg P(s) = a\deg p_1 + b \deg p_2 + (m-a)\deg q_1+(m-b)\deg q_2 \le 2mn.
    \]
    Again, we have the same situation where $P(S) \subseteq k[t]_{\deg \le 2mn}$, but $|S| = \frac{1}{2}(m+2)(m+1)$ while $\dim(k[t]_{\deg\le 2mn}) = 2mn+1$. Since dimension scales on the order of $m$ while $|S|$ scales on the order of $m^2$, for large enough $m$ we get $|S| > \dim (k[t]_{\deg \le 2mn})$ and thus that $S$ is $k$-linearly dependent.

    Let $C$ be any rational curve in $k^2$ given by $x=p_1(t)/q_1(t)$, $y=p_2(t)/q_2(t)$. By the linear dependence result above, find some $m$ that works to let \[F(x,y,z,w) = \sum_{i,j\in \mathcal T_F} c_{ij} x^i y^jz^{m-i}w^{m-j} \in k[x,y,z,w]\] be some nonzero polynomial such that $F(p_1(t),p_2(t),q_1(t),q_2(t))=0$ for all $t \in k$. Define a new nonzero polynomial $G \in k[x,y]$ simply by dropping the $z$ and $w$ variables from $F$, replacing them with $1$ wherever they appear. Consider that
    \[
    G\left(\frac{p_1(t)}{q_1(t)}, \frac{p_2(t)}{q_2(t)}\right) = \sum_{i,j\in \mathcal T_F} c_{ij} \frac{p_1(t)^i}{q_1(t)^i}\frac{p_2(t)^j}{q_2(t)^j}.
    \]
    We wish to show that $C \subseteq \mathbf V(G)$, i.e. that for all $t \in k$ where $q_1(t)\ne 0$ and $q_2(t)\ne 0$, the expression above evaluates to 0. For any $t \in k$, we can evaluate
    \begin{align*}
        q_1(t)^m q_2(t)^m G\left(\frac{p_1(t)}{q_1(t)},\frac{p_2(t)}{q_2(t)} \right) &= \sum_{i,j\in\mathcal T_F} c_{ij} p_1(t)^i p_2(t)^j q_1(t)^{m-i}q_2(t)^{m-j} \\
        &= F(p_1(t),p_2(t),q_1(t),q_2(t))=0.
    \end{align*}
    The division step is valid since by the definition of $G$ and $F$ we know that for all terms, $i \le m$ and $j \le m$ since $i + j \le m$. In the case that $q_1(t) \ne 0$ and $q_2(t) \ne 0$, by the fact that fields are integral domains we know that $G\left(\frac{p_1(t)}{q_1(t)},\frac{p_2(t)}{q_2(t)}\right) = 0$. Thus, we have shown our desired result $C \subseteq \mathbf V(G)$ from the fact that $G$ is zero whenever both denominators of rational functions are nonzero!
\end{proof}
\end{prob}
\bigskip 

I love writing the same proof 3 times with slight differences.


\end{flushleft}
\end{document}